% Baccalauréat malgache, série D, 2010 : énoncé.
\TCDBacAttribution{D}{2010}
\providecommand{\Courbe}{(\mathscr{C})}
\providecommand{\D}{(\mathscr{D})}
\providecommand{\Plan}{(\mathscr{P})}
\providecommand{\Ouv}{(O ; \vec{u} , \vec{v})}
%\providecommand{\Oij}{(O ; \vec{\imath} , \vec{\jmath})}
\providecommand{\cm}{~\mathrm{cm}}
\providecommand{\dx}{~\mathrm{d}x}
\providecommand{\Arg}[1]{\mathrm{Arg}\left( #1 \right)}

\noindent \underline{\textsc{Exercice I}} (5 points) \newline
      Le plan complexe $\Plan$ est muni d'un rep\`ere orthonorm\'e directe $\mathfrak{R} = \Ouv$ d'unit\'e $2\cm$.
      \begin{enumerate}
            \item On consid\`ere les points $S$, $R$ et $D$ d'affixes respectives $(-4)$, $(1-i)$ et $(3i)$.
                  Repr\'esenter, dans $\mathfrak{R}$, l'image du triangle $SRD$, par la similitude plane directe 
                  de centre $R$, d'angle $-\frac{\pi}{2}$ et de rapport $\frac{1}{2}$.
                  \pnt{0,5}
            \item Soit $\psi$ le nombre complexe d\'efini par $\psi = \frac{z-3i}{z+4}$.
                  \begin{enumerate}
                        \item Interpr\'eter g\'eom\'etriquement $\psi$.
                              \point{0,25}{0,25}
                        \item $M$ \'etant le point d'affixe $z$, pr\'eciser la nature de chacun des ensembles $(\Delta)$
                              et $(\Sigma)$ d\'efinis respectivement par : 
                              $
                                    (\Delta) = \left\lbrace M \in \Plan \mid |\psi| = 1 \right\rbrace
                              $
                              et 
                              $ 
                                    (\Sigma) = \left\lbrace M \in \Plan \mid \Arg{\psi} = \frac{\pi}{2} \right\rbrace
                                    .
                              $
                              \Points{0,5}{0,5}{0,25}
                              On justifiera que $(\Sigma)$ ne passe pas par l'origine des axes.
                        \item Construire, dans $\mathfrak{R}$, $(\Delta)$ et $(\Sigma)$.
                              \pnt{1}
                  \end{enumerate}
            \item On donne le polyn\^ome $Q$, de la variable complexe $z$, d\'efini par : 
                  $Q(z) = z^4 - 3i z^3 + 8iz + 24$.
                  \begin{enumerate}
                        \item Calculer $Q(3i)$.
                              \pnt{0,5}
                        \item Ecrire sous forme trigonom\'etrique les solutions, dans $\mathbb{C}$, de l'\'equation $Q(z)=0$.
                              \pnts{1,25}
                  \end{enumerate}
      \end{enumerate}
%%%%%%%%%%%%%%%%%%%%%%%%%%%%%%%%%%%%%%%%%%%%%%%%%%%%%%%%%%%%%%%%%%%%%%%%%%%%%%%%%%%%%%%%%%%%%%%
%%%%%%%%%%%%%%%%%%%                   %%%%%%%%%%%%%%%%%%%                   %%%%%%%%%%%%%%%%%%%
%%%%%%%%%%%%%%%%%%%%%%%%%%%%%%%%%%%%%%%%%%%%%%%%%%%%%%%%%%%%%%%%%%%%%%%%%%%%%%%%%%%%%%%%%%%%%%%
\noindent \underline{\textsc{Exercice II}} (5 points) 
      \begin{enumerate}
            \item A l'issue d'une exp\'erience al\'eatoire, on d\'efinit une variable al\'eatoire $X$ par le tableau suivant :
                  $$
                  \begin{array}{|c|c|c|c|c|}
                  \hline 
                  x & -1 & 2 & 3 & 4 \\ 
                  \hline 
                  p(X = x) & \frac{1}{16} & \frac{5}{16} & \frac{3}{16} & \frac{7}{16} \\ 
                  \hline 
                  \end{array} 
                  $$
                  \begin{enumerate}
                        \item D\'efinir la fonction de r\'epartition $F$ de $X$. \pnt{1}
                        \item Repr\'esenter graphiquement $F$ dans un rep\`ere orthogonal $(O;\vec{\imath},\vec{\jmath})$
                              tel que $\| \vec{\imath} \| = 2 \cm$ et $\| \vec{\jmath} \| = 8 \cm$.
                              \pnt{1}
                        \item Calculer l'esp\'erance math\'ematique et la variance de $X$.
                              \point{0,5}{0,5}
                  \end{enumerate}
            \item Une soubique contient des poussins jaunes et des poussins gris.
                  Un enfant prend au hasard et simultan\'ement deux poussins de la soubique.
                  Il y a 3 fois plus de chances, pour lui, de prendre deux gris que deux jaunes.
                  La probabilit\'e de prendre deux poussins de couleurs diff\'erentes est 
                  le double de celle de prendre deux poussins gris.
                  \begin{enumerate}
                        \item Prouver que la probabilit\'e pour l'enfant de prendre deux poussins gris est \'egale \`a $0,3$.
                              \pnts{1,5}
                        \item D\'eterminer la probabilit\'e pour l'enfant de prendre au moins un poussin gris. 
                              \pnt{0,5}
                  \end{enumerate}
      \end{enumerate}
%%%%%%%%%%%%%%%%%%%%%%%%%%%%%%%%%%%%%%%%%%%%%%%%%%%%%%%%%%%%%%%%%%%%%%%%%%%%%%%%%%%%%%%%%%%%%%%
%%%%%%%%%%%%%%%%%%%                   %%%%%%%%%%%%%%%%%%%                   %%%%%%%%%%%%%%%%%%%
%%%%%%%%%%%%%%%%%%%%%%%%%%%%%%%%%%%%%%%%%%%%%%%%%%%%%%%%%%%%%%%%%%%%%%%%%%%%%%%%%%%%%%%%%%%%%%%
\noindent \underline{\textsc{Probleme}} \newline
      On consid\`ere la fonction num\'erique $f$ d\'efinie par : 
      \[
            \begin{multiform}
                  f(0) = 0 \\
                  f(x) = - x \ln (x) - (1-x) \ln (1-x) \text{ si } 0 < x < 1 \\
                  f(1) = 0 \;.
            \end{multiform}
      \]
      \begin{enumerate}
            \item Pr\'eciser l'ensemble de d\'efinition de $f$.
                  \pnt{0,25}
            \item Prouver que $f$ est continue en $0$ et en $1$.
                  \pnt{1}
            \item Justifier que l'on a : 
                  \[
                        \lim_{x\rightarrow 0^+} \frac{f(x) - f(0)}{x - 0} = +\infty
                        \text{ et } 
                        \lim_{x\rightarrow 1^-} \frac{f(x) - f(1)}{x - 1} = -\infty \;.
                  \]
                  \point{0,5}{0,5}
                  Que peut-on d\'eduire pour la d\'erivabilit\'e de $f$ en $0$ et en $1$, 
                  puis pour les tangentes \`a la courbe repr\'esentative de $f$ aux points 
                  $O(0;0)$ et $P(1;0)$\,?
                  \POINT{0,25}{0,25}{0,25}{0,25}
            \item \begin{enumerate}
                        \item Calculer $f'(x)$ pour tout r\'eel $x$ de l'intervalle $]0;1[$.
                              \pnt{0,75}
                        \item R\'esoudre, dans $]0;1[$, l'in\'equation : $\ln (1-x) - \ln (x) \leq 0$.
                              \pnt{0,5}
                        \item Dresser le tableau de variation de $f$.
                              \pnt{0,75}
                  \end{enumerate}
            \item Tracer la courbe repr\'esentative de $f$ et ses tangentes, dans le m\^eme rep\`ere orthonorm\'e 
                  \Oij{} d'unit\'e $10\cm$.
                  \pnts{1,75}
                  
                  (On prend $\ln (2) \simeq 0,7$)
            \item Soient $u$ et $v$ les fonctions num\'eriques d\'efinies respectivement sur $]0;1[$ par : 
                  \[
                        u(x) = \frac{x^2}{2} \ln (x) - \frac{x^2}{4}
                        \text{ et }
                        v(x) = \frac{(x-1)^2}{2} \ln (1-x) - \frac{x^2 - 2x}{4} \;.
                  \]
                  \begin{enumerate}
                        \item Calculer $u'(x)$ et $v'(x)$.
                              \point{0,5}{0,5}
                        \item D\'eterminer la primitive sur $]0;1[$ de $f$ qui prend la valeur $-\frac{3}{4}$ en $\frac{1}{2}$.
                              \pnts{2}
                  \end{enumerate}
      \end{enumerate}











